Shear stress in timber occurs when forces act parallel to the grain direction, causing one layer of fibers to slide over another along a plane parallel to the grain. This type of failure is particularly critical in notched beams, connections, and lap joints. \\[4mm]
\noindent Wood is surprisingly weak when forces try to slide one layer of fibers over another along the grain, because only the weak lignin glue between fibers is resisting it — making shear its sneakiest failure mode. Unlike bending failures that show visible signs, shear failure splits the beam suddenly and horizontally with absolutely no warning. This is why engineers must always specifically check shear at vulnerable points like notches and connections, because that is exactly where timber will silently choose to fail.